\section{Исследование и построение решения задачи}
\label{sec:Section3} \index{Section3}

\textbf{Унифицированная вендоро-независимая модель}. Модель должна позволять представлять правила различных платформ (классические ACL Cisco/Juniper/Huawei, отечественные NGFW, iptables) единообразно. На уровне атрибутов это означает поддержку как примитивных квалификаторов (адреса, порты, протоколы), так и именованных объектов и групп. На уровне действий --- базовое разделение на \emph{allow}/\emph{deny} с возможностью различения дополнительных классов действий и состояния \emph{Undecided} в смысле Diekmann для chain-семантик.

\textbf{Обработка неизвестных match-условий}. Опираясь на работу Diekmann, в систему должна быть встроена техника upper/lower closure для match-условий, которые невозможно интерпретировать на уровне модели. При этом аппроксимация должна быть информативной: вместо тотальной потери информации, для неизвестных match-условий следует использовать пятёрку из правила как контекст, остальные поля считая допускающими произвольное значение.

\textbf{Поиск интрафайервольных аномалий полного покрытия}. По образцу FIREMAN система должна оперировать состоянием $\langle A,D,F\rangle$ или его эквивалентом и обнаруживать классы shadowing, redundancy, generalization, correlation, а также irrelevance. Union-аномалии могут быть обнаружены, но представляются предупреждением, а не ошибкой, ввиду указанной выше специфики управления правилами в Zero-Trust-инфраструктурах.

\textbf{Опора на методологию статического post-change аудита}. Разрабатываемая система предназначена для регулярной проверки уже эксплуатируемых конфигураций, а не для анализа предполагаемых изменений. Это определяет приоритеты: точность и полнота обнаружения аномалий важнее, чем интерактивная производительность; объём отчёта должен быть достаточным для администратора, но не разрозненным; результат должен быть привязан к конкретным правилам исходной конфигурации.

\textbf{Двухфазная архитектура нормализатор + анализатор}. Парсинг и нормализация исходных конфигураций отделяются от ядра анализа в самостоятельный модуль. Это обеспечивает расширяемость на новые форматы и позволяет вводить поддержку отечественных платформ инкрементно, без переработки ядра.

\textbf{Возможность расширения}. Архитектура системы должна оставлять место для функций, не входящих в основной набор: change-impact-анализ, log-mining, reachability-анализ, интеграция с системами учёта правил.

С учётом перечисленных требований основными ориентирами для архитектуры разрабатываемого инструмента являются работы Yuan и др.~\cite{yuan2006fireman} и Diekmann и др.~\cite{diekmann2018verified}. От FIREMAN заимствуются формальная модель состояний $\langle A,D,F\rangle$ при анализе ACL, классификация аномалий и структура отчёта. От Diekmann --- двухфазная архитектура нормализатор + анализатор, тернарная логика для неизвестных match-условий, идея сервисных матриц как формы сводного отчёта. Дополнительно используется унификация представления правил, описанная в обзоре Ramesh и др.~\cite{ramesh2024exploring} (отношения IM/EM/PM/CD/C). Совокупность этих идей служит основанием для разрабатываемой в последующих разделах унифицированной вендоро-независимой модели правил и алгоритмов анализа на её основе.
